% Build: pdflatex -> biber -> pdflatex -> pdflatex (or just latexmk -pdf)
% Add the `handout` class option to collapse all \pause overlays.
\PassOptionsToPackage{prologue,dvipsnames}{xcolor}
\PassOptionsToPackage{colorlinks=true,linkcolor=.,urlcolor=mLightBrown,citecolor=mLightBrown,hyperfootnotes=true}{hyperref}
\documentclass[aspectratio=169,10pt]{beamer}
% ---------- Theme: Moloch (maintained fork of Metropolis) ----------
\usetheme{moloch}
\molochset{
    progressbar=frametitle,   % thin progress bar under each frame title
    numbering=fraction,       % "3 / 32" in the footer
    block=fill,               % filled blocks
    sectionpage=progressbar,  % section title slides with a progress bar
    subsectionpage=none,
}
\setbeamertemplate{theorems}[normal font]   % upright text in theorem-like blocks

% ---------- Packages ----------
\usepackage[english]{babel}
\usepackage[utf8]{inputenc}
\usepackage[T1]{fontenc}
\usepackage{csquotes}
\usepackage{etoolbox}
\usepackage{float}
\usepackage{multirow}
\usepackage{array}
\usepackage{enumerate}
\usepackage{booktabs}
\usepackage{makecell}
\usepackage{adjustbox}
\usepackage{setspace}
\usepackage{subcaption}
\usepackage{pgfpages}
\usepackage{appendixnumberbeamer}   % frames after \appendix don't count toward the total

\usepackage{amsmath,amsfonts,amssymb,amsthm,mathtools,mathrsfs}

% ---------- Bibliography ----------
\usepackage[style=alphabetic,maxcitenames=1,backend=biber]{biblatex}
\renewcommand*{\bibfont}{\scriptsize}
\AtBeginBibliography{\let\emph\textit}  % plain italics in the bibliography
\setcounter{biburlnumpenalty}{9000}
\setcounter{biburlucpenalty}{9000}
\setcounter{biburllcpenalty}{9000}

% ---------- Colors ----------
% mDarkTeal (#23373B) and mLightBrown (#EB811B) are defined by Moloch.
\definecolor{cGoal}{HTML}{EB811B}
\definecolor{cProblem}{HTML}{C0392B}
\definecolor{cIntuition}{HTML}{7B5EA7}
\definecolor{cRemark}{HTML}{0E8388}
\definecolor{cPrev}{HTML}{7F8C8D}
\definecolor{cAssumption}{HTML}{2874A6}
\definecolor{cObserve}{HTML}{AD1457}
\definecolor{cExample}{HTML}{2E7D32}

% Colored, italic emphasis; \emph<2->{...} colors immediately and italicizes on later overlays
\renewcommand<>{\emph}[1]{{\color{mLightBrown}\only#2{\itshape}#1}}

% Footnotes numbered per frame
\let\oldfootnote\footnote
\renewcommand\footnote[1][]{\oldfootnote[frame,#1]}

% ---------- Colored block environments ----------
% \styleblock{color} recolors the *next* block/theorem-like environment in the
% current group, e.g. {\styleblock{cExample}\begin{block}{Title} ... \end{block}}
\newcommand{\styleblock}[1]{%
    \setbeamercolor{block title}{fg=white,bg=#1}%
    \setbeamercolor{block body}{fg=black,bg=#1!10!white}%
}
\newtheorem*{goal}{Goal}
\AtBeginEnvironment{goal}{\styleblock{cGoal}}
\newtheorem*{intuition}{Intuition}
\AtBeginEnvironment{intuition}{\styleblock{cIntuition}}
\newtheorem*{remark}{Remark}
\AtBeginEnvironment{remark}{\styleblock{cRemark}}
\newtheorem*{prev}{As previously seen}
\AtBeginEnvironment{prev}{\styleblock{cPrev}}
\newtheorem*{assumption}{Assumption}
\AtBeginEnvironment{assumption}{\styleblock{cAssumption}}
\newtheorem*{observe}{Observe}
\AtBeginEnvironment{observe}{\styleblock{cObserve}}
\AtBeginEnvironment{problem}{\styleblock{cProblem}}   % `problem` is predefined by beamer

% ---------- TodoNotes and inline notes in fancy boxes ----------
\usepackage{todonotes}
\usepackage{tcolorbox}

% ---------- quiver style ----------
\usepackage{tikz-cd}
% `calc` is necessary to draw curved arrows.
\usetikzlibrary{calc}
% `pathmorphing` is necessary to draw squiggly arrows.
\usetikzlibrary{decorations.pathmorphing}

% A TikZ style for curved arrows of a fixed height, due to AndréC.
\tikzset{curve/.style={settings={#1},to path={(\tikztostart)
                    .. controls ($(\tikztostart)!\pv{pos}!(\tikztotarget)!\pv{height}!270:(\tikztotarget)$)
                    and ($(\tikztostart)!1-\pv{pos}!(\tikztotarget)!\pv{height}!270:(\tikztotarget)$)
                    .. (\tikztotarget)\tikztonodes}},
    settings/.code={\tikzset{quiver/.cd,#1}
            \def\pv##1{\pgfkeysvalueof{/tikz/quiver/##1}}},
    quiver/.cd,pos/.initial=0.35,height/.initial=0}

% TikZ arrowhead/tail styles.
\tikzset{tail reversed/.code={\pgfsetarrowsstart{tikzcd to}}}
\tikzset{2tail/.code={\pgfsetarrowsstart{Implies[reversed]}}}
\tikzset{2tail reversed/.code={\pgfsetarrowsstart{Implies}}}
% TikZ arrow styles.
\tikzset{no body/.style={/tikz/dash pattern=on 0 off 1mm}}

% ---------- Math macros ----------
\newcommand{\at}[3]{\left.#1\right\vert_{#2}^{#3}}
\newcommand\quotient[2]{
    \mathchoice
    {% \displaystyle
        \text{\raise1ex\hbox{$#1$}\Big/\lower1ex\hbox{$#2$}}%
    }
    {% \textstyle
        #1\,/\,#2
    }
    {% \scriptstyle
        #1\,/\,#2
    }
    {% \scriptscriptstyle
        #1\,/\,#2
    }
}

\newcommand{\identity}{\mathrm{id}}
\newcommand{\Homomorphism}{\mathrm{Hom}}
\newcommand{\Morphism}{\mathrm{Mor}}
\newcommand{\Object}{\mathrm{Ob}}
\newcommand{\True}{\textsf{True}}
\newcommand{\T}{\textsf{T}}
\newcommand{\False}{\textsf{False}}
\newcommand{\F}{\textsf{F}}
\newcommand{\act}{\rotatebox[origin=c]{-180}{\(\,\circlearrowright\,\)}}

\usepackage{physics}
\usepackage{complexity}

\DeclareMathOperator{\im}{Im}
\DeclareMathOperator{\sgn}{sgn}
\DeclareMathOperator{\Int}{Int}
\DeclareMathOperator{\diag}{diag}
\DeclareMathOperator{\dom}{dom}
\DeclareMathOperator{\OPT}{\textsf{OPT}}
\DeclareMathOperator*{\argmax}{arg\,max}
\DeclareMathOperator*{\argmin}{arg\,min}
\DeclareMathOperator{\Var}{Var}
\DeclareMathOperator{\Cov}{Cov}

\let\implies\Rightarrow
\let\impliedby\Leftarrow
\let\iff\Leftrightarrow

\usepackage{stmaryrd} % for \lightning
\newcommand\conta{\scalebox{1.1}{\(\lightning\)}}

\usepackage{bm}
\usepackage{bbm}

% ---------- Figure support ----------
\usepackage{import}
\usepackage{xifthen}
\usepackage{pdfpages}
\usepackage{transparent}
\newcommand{\incfig}[2][\columnwidth]{%
    \def\svgwidth{#1}
    \import{Figures/}{#2.pdf_tex}
}

\makeatletter
\let\@@magyar@captionfix\relax
\makeatother

\newenvironment<>{proofs}[1][\proofname]{%
    \par
    \def\insertproofname{#1.}%
    \usebeamertemplate{proof begin}#2}
{\usebeamertemplate{proof end}}

% logo of my university
\titlegraphic{\hfill\includegraphics[height=1.1cm]{~/Developer/LaTeX-Template/Beamer/Figures/logo/logo2.png}}


\title{Beamer Template}
\subtitle{A Demo of the Modern Style}
\date{\today}
\author{Pingbang Hu}
\institute{University of Illinois Urbana-Champaign}

\addbibresource{reference.bib}

\begin{document}

\maketitle

\begin{frame}{Outline}
    \tableofcontents[hideallsubsections]
\end{frame}

% =====================================================================
\section{Useful Environment}
% =====================================================================

\begin{frame}{\texttt{itemize} and \texttt{enumerate}}
    This is a demonstration of various environments~\cite{HuLaTeXTemplate}. The first is \texttt{itemize}:
    \begin{itemize}
        \item Item 1
        \item Item 2
        \item etc.
    \end{itemize}
    \pause

    \medskip
    Usually I'll use \texttt{enumerate} as well, revealed step by step with \texttt{\textbackslash pause}:
    \begin{enumerate}
        \item Item 1 \pause
        \item Item 2 \pause
        \item etc.
    \end{enumerate}
    \pause

    \medskip
    Use \texttt{\textbackslash emph} to \emph{highlight} a phrase, and \texttt{\textbackslash footnote} for asides.\footnote{Footnotes are numbered per frame.}
\end{frame}

\begin{frame}{Definition, Examples, etc.}
    All the common environments are available, e.g.,
    \begin{itemize}
        \item \emph{Theorem}-related: \texttt{lemma}, \texttt{theorem}, \texttt{corollary}.
        \item \emph{Common} environments: \texttt{definition}, \texttt{example}, \texttt{problem}.
        \item \emph{Others}: \texttt{goal}, \texttt{prev}, \texttt{observe}, \texttt{assumption}, \texttt{remark}, \texttt{intuition}.
    \end{itemize}
    \pause
    \begin{remark}
        Each custom environment has its own color; the name is the block title, so they are unnumbered.
    \end{remark}
    \pause
    \begin{goal}
        A \texttt{goal} environment.
    \end{goal}
    \pause
    \begin{prev}
        A \texttt{prev} environment.
    \end{prev}
\end{frame}

\begin{frame}{More Blocks}
    \begin{columns}[t]
        \begin{column}{0.48\textwidth}
            \begin{problem}
                A \texttt{problem} environment.
            \end{problem}
            \begin{assumption}
                An \texttt{assumption} environment.
            \end{assumption}
            \begin{observe}
                An \texttt{observe} environment.
            \end{observe}
        \end{column}
        \begin{column}{0.48\textwidth}
            \begin{intuition}[Optional title]
                An \texttt{intuition} environment.
            \end{intuition}
            \begin{example}
                A plain \texttt{example} environment.
            \end{example}
            {\styleblock{cExample}%
                \begin{block}{Ad-hoc color}
                    Wrap a \texttt{block} with \texttt{\textbackslash styleblock\{color\}} for a one-off color.
                \end{block}}
        \end{column}
    \end{columns}
\end{frame}

% =====================================================================
\section{Math and Figures}
% =====================================================================

\begin{frame}{Math}
    The usual macros are available, e.g., \(\argmin_{\theta} \Var[f_\theta(x)]\), \(\dd{x}\), \(\norm{\bm{v}}\), and \(\P \neq \NP\).
    \pause
    \begin{theorem}[A statement]
        For every \(\epsilon > 0\) there exists \(\delta > 0\) such that \(\abs{x - y} < \delta \implies \abs{f(x) - f(y)} < \epsilon\).
    \end{theorem}
    \pause
    \begin{proof}
        Left as an exercise. Use \texttt{proofs} for a proof that spans overlays.
    \end{proof}
\end{frame}

\begin{frame}[fragile]{Figures}
    Use \texttt{\textbackslash includegraphics} as usual, or \texttt{\textbackslash incfig\{name\}} for Inkscape figures exported as \texttt{Figures/name.pdf\_tex}.
    \begin{center}
        \begin{tikzcd}
            A \arrow[r, "f"] \arrow[dr, "g \circ f"'] & B \arrow[d, "g"] \\
            & C
        \end{tikzcd}
    \end{center}
\end{frame}

\begin{frame}[standout]
    Thanks! Ask \textbf{\emph{anything}} you want!
\end{frame}

\appendix

\begin{frame}[allowframebreaks]{References}
    \printbibliography[heading=none]
\end{frame}

\end{document}
